\subsection{CONSEQUENCES OF DISTRIBUTIVITY}

Distributivity of multiplication with respect to addition allows us to apply Proposition 1 of §3, no. 4, which gives

\[
\prod_ {i = 1} ^ {n} \left(\sum_ {\lambda \in L _ {i}} x _ {i, \lambda}\right) = \sum_ {\alpha_ {1}, \dots , \alpha_ {n}} \prod_ {i = 1} ^ {n} x _ {i, \alpha_ {i}}\tag{13}
\]

where the sum extends over all sequences \((\alpha_{1},\ldots,\alpha_{n})\) belonging to \(L_{1}\times\cdots\times L_{n}\) and for \(i=1,\ldots,n\) the family \((x_{i,\lambda})_{\lambda\in L^{i}}\) of elements of the ring A is of finite support.

PROPOSITION 1. Let A be a commutative ring and \((x_{\lambda})_{\lambda \in L}\) a finite family of elements of A. For every family of positive integers \(\beta = (\beta_{\lambda})_{\lambda \in L}\) , let \(|\beta| = \sum_{\lambda \in L} \beta_{\lambda}\) . Then

\[
\left(\sum_ {\lambda \in L} x _ {\lambda}\right) ^ {n} = \sum_ {| \beta | = n} \frac {n !}{\prod_ {\lambda \in L} \beta_ {\lambda} !} \prod_ {\lambda \in L} x _ {\lambda} ^ {\beta_ {\lambda}}.\tag{14}
\]

We apply formula (13) with \(\mathbf{L}_i = \mathbf{L}\) and \(x_{i,\lambda} = x_{\lambda}\) for \(1 \leqslant i \leqslant n\) . Then

\[
\left(\sum_ {\lambda \in L} x _ {\lambda}\right) ^ {n} = \sum_ {\alpha_ {1}, \dots , \alpha_ {n}} x _ {\alpha_ {1}} \dots x _ {\alpha_ {n}},\tag{15}
\]

the sum extending over all sequences \(\alpha = (\alpha_{1},\ldots ,\alpha_{n})\in \mathbf{L}^{n}\)

Let \(\alpha\) be in \(\mathbf{L}^n\) ; for all \(\lambda \in \mathbf{L}\) let \(\mathrm{U}_{\lambda}^{\alpha}\) denote the set of integers \(i\) such that \(1 \leqslant i \leqslant n\) and \(\alpha_i = \lambda\) and let \(\Phi(\alpha) = (\mathrm{U}_{\lambda}^{\alpha})_{\lambda \in \mathbf{L}}\) . It is immediate that \(\Phi\) is a bijection of \(\mathbf{L}^n\) onto the set of partitions of \(\{1, 2, \ldots, n\}\) indexed by \(\mathbf{L}\) . For all \(\beta \in \mathbf{N}^{\mathbf{L}}\) such that \(|\beta| = n\) , let \(\mathbf{L}_{\beta}^n\) denote the set of \(\alpha \in \mathbf{L}^n\) such that \(\operatorname{Card} \mathrm{U}_{\beta}^{\alpha} = \beta_{\lambda}\) for all \(\lambda \in \mathbf{L}\) . It follows that the family \((\mathbf{L}_{\beta}^n)_{\beta | = n}\) is a partition of \(\mathbf{L}^n\) and that

\[
\operatorname{Card} \mathrm{L} _ {\beta} ^ {n} = \frac {n !}{\prod_ {\lambda \in \mathrm{L}} \beta_ {\lambda} !}
\]

(§ 5, no. 5).

Finally, for \(\alpha \in \mathbf{L}_{\beta}^{n}\) ,

\[
x _ {\alpha_ {1}} \dots x _ {\alpha_ {n}} = \prod_ {\lambda \in L} \prod_ {i \in U _ {\lambda} ^ {\alpha}} x _ {\alpha_ {i}} = \prod_ {\lambda \in L} \prod_ {i \in U _ {\lambda} ^ {\alpha}} x _ {\lambda} = \prod_ {\lambda \in L} x _ {\lambda},
\]

whence

\[
\begin{array}{r l} \sum_ {(\alpha_ {1}, \dots , \alpha_ {n})} x _ {\alpha_ {1}} \dots x _ {\alpha_ {n}} & = \sum_ {| \beta | = n} \sum_ {\alpha \in L _ {\beta} ^ {n}} x _ {\alpha_ {1}} \dots x _ {\alpha_ {n}} \\ & = \sum_ {| \beta | = n} \sum_ {\alpha \in L _ {\beta} ^ {n}} \prod_ {\lambda \in L _ {\lambda}} x _ {\lambda} ^ {\beta_ {\lambda}} \\ & = \sum_ {| \beta | = n} \frac {n !}{\prod_ {\lambda \in L} \beta_ {\lambda} !} x _ {\lambda} ^ {\beta_ {\lambda}} \end{array}
\]

and formula (14) thus follows from (15).

COROLLARY 1 (binomial formula). Let \(x\) and \(y\) be two elements of a commutative ring A. Then:

\[
(x + y) ^ {n} = \sum_ {p = 0} ^ {n} \binom {n} {p} x ^ {p} y ^ {n - p}.
\]

Formula (14) applied to \(\mathbf{L} = \{1,2\}\) , \(x_{1} = x\) and \(x_{2} = y\) gives

\[
(x + y) ^ {n} = \sum_ {p + q = n} \frac {n !}{p ! q !} x ^ {p} y ^ {q},
\]

the sum extending over ordered pairs of positive integers \(p, q\) with \(p + q = n\) . The binomial formula follows immediately from this (Set Theory, III, § 5, no. 8).

COROLLARY 2. Let A be a commutative ring, X a set, \(\mathbf{u} = (u_x)_{x \in \mathbb{X}}\) and \(\mathbf{v} = (v_x)_{x \in \mathbb{X}}\) two families of elements of A. Let \(\mathbf{u} + \mathbf{v}\) denote the family \((u_x + v_x)_{x \in \mathbb{X}}\) . For all \(\lambda \in \mathbf{N}^{(\mathbb{X})}\) we write \(\lambda! = \prod_{x \in \mathbb{X}} \lambda(x)!\) . Then for all \(\alpha \in \mathbf{N}^{(\mathbb{X})}\) , in the notation of § 7, no. 8,

\[
(\mathbf {u} + \mathbf {v}) ^ {\alpha} = \sum_ {\beta + \gamma = \alpha} \frac {\alpha !}{\beta ! \gamma !} \mathbf {u} ^ {\beta} \mathbf {v} ^ {\gamma}.
\]

For \(x\in \mathbf{X}\)

\[
(u _ {x} + v _ {x}) ^ {\alpha (x)} = \sum_ {m + n = \alpha (x)} \frac {\alpha (x) !}{m ! n !} u _ {x} ^ {m} v _ {x} ^ {n}
\]

by Corollary 1. Taking the product of these equations for \(x \in \mathbf{X}\) and using (13), we obtain the corollary.

PROPOSITION 2. Let A be a ring, \(x_{1}, \ldots, x_{n}\) elements of A and I = \{1, 2, \ldots, n\}. For H ⊂ I, we write \(x_{\mathrm{H}} = \sum_{i \in \mathrm{H}} x_{i}\) . Then

\[
(- 1) ^ {n} \sum_ {\sigma \in \mathfrak {S} _ {n}} x _ {\sigma (1)} \dots x _ {\sigma (n)} = \sum_ {H \subset I} (- 1) ^ {\operatorname{Card} H} (x _ {H}) ^ {n}.\tag{16}
\]

In particular, if A is commutative,

\[
(- 1) ^ {n} n! x _ {1} x _ {2} \dots x _ {n} = \sum_ {\mathrm{H} \in \mathrm{I}} (- 1) ^ {\text { Card   H }} (x _ {\mathrm{H}}) ^ {n}.
\]

Let C be the set of mappings of I into \(\{0,1\}\) . If each \(H \subset I\) is mapped to its characteristic function, a bijection is obtained of \(\mathfrak{P}(I)\) onto C. The right hand side of (16) is thus equal to:

\[
\begin{array}{r l} \sum_ {a \in \mathbb {C}} (- 1) ^ {a (1) + \dots + a (n)} & \left(\sum_ {i \in \mathbb {I}} a (i) x _ {i}\right) ^ {n} \\ & = \sum_ {a \in \mathbb {C}} (- 1) ^ {a (1) + \dots + a (n)} \sum_ {(i _ {1}, \ldots , i _ {n}) \in \mathbb {I} ^ {n}} a (i _ {1}) \ldots a (i _ {n}) x _ {i _ {1}} \ldots x _ {i _ {n}} \\ & = \sum_ {(i _ {1}, \ldots , i _ {n}) \in \mathbb {I} ^ {n}} c _ {i _ {1} \ldots i _ {n}} x _ {i _ {1}} \ldots x _ {i _ {n}} \end{array}
\]

where

\[
c _ {i _ {1} \dots i _ {n}} = \sum_ {a \in \mathbb {C}} (- 1) ^ {a (1) + \dots + a (n)} a (i _ {1}) \dots a (i _ {n}).
\]

\begin{enumerate}
\def\labelenumi{(\arabic{enumi})}
\tightlist
\item
  Suppose that \((i_{1},\ldots,i_{n})\) is not a permutation of I. There exists a \(j\in I\) distinct from \(i_{1},\ldots,i_{n}\) . Let \(C'\) be the set of \(a\in C\) such that \(a(j)=0\) . For all \(a\in C'\) , let \(a^{*}\) be the sum of a and the characteristic function of \(\{j\}\) . Then \(a^{*}(1)+\cdots+a^{*}(n)=a(1)+\cdots+a(n)+1\) and hence
\end{enumerate}

\[
\begin{array}{l} c _ {i _ {1} \dots i _ {n}} = \sum_ {a \in C ^ {\prime}} (- 1) ^ {a (1) + \dots + a (n)} a (i _ {1}) \dots a (i _ {n}) + (- 1) ^ {a ^ {*} (1) + \dots + a ^ {*} (n)} a ^ {*} (i _ {1}) \dots a ^ {*} (i _ {n}) \\ = \sum_ {a \in C ^ {\prime}} ((- 1) ^ {a (1) + \dots + (n)} + (- 1) ^ {a (1) + \dots + a (n) + 1}) a (i _ {1}) \dots a (i _ {n}) = 0. \end{array}
\]

\begin{enumerate}
\def\labelenumi{(\arabic{enumi})}
\setcounter{enumi}{1}
\tightlist
\item
  Suppose that there exists \(\sigma \in \mathfrak{S}_n\) such that \(i_1 = \sigma(1), \ldots, i_n = \sigma(n)\) . Then \(a(i_1) \ldots a(i_n) = 0\) unless \(a\) only takes the value 1. Thus \(c_{i_1 \ldots i_n} = (-1)^n\) .
\end{enumerate}

